\par\medskip
\begin{center}\hypertarget{native-164-article118}{}\pdfbookmark[2]{118.}{article-118}118.\end{center}
If the distance of the celestial body from the Earth corresponding to a given observation is already approximately known, the observation can be corrected for the effect of \emph{aberration} in several ways, based on the different methods given in Article 71. Let $t$ be the actual time of observation. Let $\theta$ be the interval in which light travels from the celestial body to the Earth, found by multiplying $493^{\mathrm{s}}$ by the distance. Let $l$ be the observed position and $l'$ the same position reduced to time $t+\theta$ by means of the daily geocentric motion. Let \SourceLeaf{165}$l''$ be the position $l$ corrected for the part of aberration that planets have in common with the fixed stars. Let $L$ be the true position of the Earth corresponding to time $t$ (that is, the tabular position increased by $20''{,}25$). Finally, let ${}'L$ be the true position of the Earth corresponding to time $t-\theta$. We then have the following.
\par I. $l$ is the true position of the celestial body as seen from ${}'L$ at time $t-\theta$
\par II. $l'$ is the true position of the celestial body as seen from $L$ at time $t$
\par III. $l''$ is the true position of the celestial body as seen from $L$ at time $t-\theta$.
\par\noindent Method I thus retains the observed position as it stands and substitutes the fictitious time $t-\theta$ for the actual time, computing the Earth's position for that same time. Method II changes only the observation, but the change requires the daily motion as well as the distance. Method III applies a correction to the observation that is independent of the distance and substitutes the fictitious time $t-\theta$ for the actual time, while retaining the Earth's position corresponding to the actual time. The first of these methods is by far the most convenient whenever the distance is already known well enough for the time correction $\theta$ to be computed with sufficient precision.

If this distance is still entirely unknown, however, none of these methods can be applied directly. The first gives the geocentric position of the celestial body, but the time and the position of the Earth, which depend on the unknown distance, are missing. In the second we have the latter two but lack the geocentric position. In the third we have the geocentric position of the celestial body and the position of the Earth, but lack the time to be used with them.

What, then, should we do in our problem if, in such a case, the solution is required to be exact with respect to aberration as well? The simplest course is certainly to determine the orbit at first with aberration neglected. Since aberration can never produce a considerable effect, the distances will certainly follow from this calculation with enough precision to correct the observations for aberration by one of the methods just described and repeat the determination of the orbit more accurately. For this purpose, the third method is far preferable to the others. In the first method, all operations that depend on the Earth's position must be started again from the beginning. In the second (which cannot even be applied unless there are enough observations to derive the daily motion), all operations that depend on the geocentric position of the celestial body must be carried out anew. In the third, by contrast, if the first calculation was already based on geocentric positions corrected for the aberration of the fixed stars, all preliminary operations that depend on the \SourceLeaf{166}Earth's position and the geocentric position of the celestial body can be carried over unchanged to the new calculation. Indeed, this method even lets us include aberration in the very first calculation if the method used to determine the orbit is such that the values of the distances are obtained before the corrected times need to be introduced into the calculation. There will then be no need for a second calculation on account of aberration, as will become clearer when we treat our problem more fully.
